%! AP Statistics — Unit 1, Lesson 1.6 — Group Activity ANSWER KEY
\documentclass[10pt]{article}
\usepackage{apstats-article}
\usepackage{apstats-key}

%=============================================================================
\begin{document}

\pageheader{Unit 1, Lesson 1.6}{Group Activity --- Answer Key}
\begin{tabularx}{\linewidth}{X X X}
Name:\;\ans{Key} & Partner:\; & Period:\;
\end{tabularx}

\vspace{0.1in}

% ── PART 1 ───────────────────────────────────────────────────────────────────
{\large\bfseries\color{navy} Part 1 \quad Model SOCS Descriptions}

\vspace{0.12in}

\noindent\textbf{Distribution A --- Ages at a Concert}

\vspace{4pt}
\ans{The distribution of ages of the 80 concert-goers is bimodal, with one
peak near 17--19 years (likely younger fans) and a second peak near 35--40
years (likely an older fan base). [S] \quad There do not appear to be any
outliers; all ages fall within the overall pattern of the distribution. [O]
\quad The center of the distribution is approximately 27 years, though a
single center value is less meaningful for a bimodal distribution. [C] \quad
The ages of these concert-goers range from 14 to 55 years, a spread of
41 years. [S]}

\vspace{0.12in}

\noindent\textbf{Distribution B --- Wait Times at a Coffee Shop}

\vspace{4pt}
\ans{The distribution of wait times for these 50 customers is strongly skewed
right, with most customers waiting between 2 and 6 minutes but a long tail
extending toward longer wait times. [S] \quad There appear to be three potential
outliers: customers who waited more than 15 minutes, with one customer waiting
22 minutes --- well above the rest of the group. [O] \quad The center of the
distribution is approximately 5 minutes, representing a typical wait time for
a customer at this coffee shop. [C] \quad The wait times range from 1 to
22 minutes, a spread of 21 minutes. [S]}

\vspace{0.12in}

\noindent\textbf{Distribution C --- Final Exam Scores}

\vspace{4pt}
\ans{The distribution of final exam scores for these 35 students is skewed
left, with most students scoring between 70 and 95 points but a tail extending
toward lower scores. [S] \quad There appear to be two potential outliers: the
two students who scored below 40 points, well below the rest of the class. [O]
\quad The center of the distribution is approximately 82 points, which
represents the typical exam score for a student in this class. [C] \quad
Scores range from 36 to 100 points, a spread of 64 points. [S]}

\vspace{0.14in}

\begin{teachernote}
Part 1 debrief: Poll students on which SOCS component was hardest. Typically:
(1) students omit spread or give only the max without the min; (2) context
drops out after the first sentence. Use the class-share to build a consensus
``strong'' response on the board before moving to Part 2.
\end{teachernote}

\vspace{0.12in}

% ── PART 2 ───────────────────────────────────────────────────────────────────
{\large\bfseries\color{navy} Part 2 \quad Critique the Sample Response}

\vspace{0.1in}

\begin{enumerate}[leftmargin=1.5em, itemsep=8pt]

  \item \ans{Problems with the sample response:\\
        (a) No context --- the variable (wait times) and population (50 coffee
        shop customers) never appear.\\
        (b) Center has no units --- ``about 5'' does not specify minutes.\\
        (c) Spread is incomplete --- ``1 to 22'' is acceptable but lacks units
        and range calculation.\\
        (d) Outliers are vague --- ``some outliers'' does not identify their
        approximate values or why they qualify.\\
        (e) Shape is correct but unsupported --- no evidence cited.}

  \item \ans{The distribution of wait times for these 50 coffee shop customers
        is strongly skewed right, with most customers waiting between 2 and
        6 minutes and a long tail extending toward longer wait times. [S]
        There appear to be three potential outliers --- customers who waited
        more than 15 minutes, including one who waited 22 minutes, which falls
        far above the rest of the group. [O] The center is approximately
        5 minutes, meaning a typical customer at this coffee shop waits about
        5 minutes. [C] Wait times range from 1 to 22 minutes, a spread of
        21 minutes. [S]}

\end{enumerate}

\end{document}
